\apendice{Plan de Proyecto Software}

\section{Introducción}

Para un buen desarrollo del proyecto se seguirá una planificación temporal para aportarnos una visión general de las tareas y plazos para facilitar el seguimiento de este trabajo y tener en cuenta complicaciones a las cuales nos debemos enfrentar. 
En este anexo se desarrollará la planificación del actual trabajo desde un punto de vista temporal y económico



\section{Planificación temporal}
GitHub fue la herramienta empleada para controlar el avance del proyecto. En repositorio se asignaron unas tareas(issues) con ciertas etiquetas(labels) para una mejor identificación, agrupándolas más tarde en hitos(milestones).
Los principales milestones son:
\begin{enumerate}
    \item \textbf{Adquisición y preprocesamiento del conjunto de datos}
En este \textit{milestone} lo que se buscaba era obtener el conjunto de datos necesarios para la implementación de la web. Los datos que se han empleado son de dos tipos, los que se obtienen mediante solicitudes api que son los de las temperaturas e índices de radiación ultravioleta, también se emplea el paquete de la INE para obtener las defunciones; y finalmente para el entrenamiento del modelo de predicción los datos han sido solicitados mediante una consulta a la delegación de la AEMET de Castilla y León.
\item \textbf{Web \textit{Shiny}}
En este apartado se agrupan las \textit{issues} del desarrollo de la plataforma. Cada \textit{issue} representa una pestaña de la web en la que se iban dejando comentarios de los avances que se realizaban o el contenido que iba a tener dicha página. A parte, se encuentra otra \textit{issue} " DAR OTROS ENFOQUES A LA APP" en la que se redactaba lo que se hablaba con en las tutorías semanales.
\item \textbf{Redacción de la memoria}
Este bloque como su propio nombre dice, se encargaba de agrupar los capitulos de la memoria que cada uno de ellos consistía en un \textit{issue}.
\item \textbf{Redaccion de los anexos}
De igual manera que el anterior, este bloque agrupaba los diferentes apéndices de los anexos.
%poner foto de la issues de cada bloque
\end{enumerate}

\section{Planificación económica}

En esta sección se analizarán los costes que supondría la realización de este trabajo en un contexto real con su respectivo personal y recursos requeridos para el desarrollo técnico de la web interactiva.
\textbf{
\subsection{{Coste personal}}}
En este proyecto solo ha habido una persona involucrada por lo que se calculará un coste de personal de un ingeniero de la salud recién graduado. Dado que no existe estimación exacta para un ingeniero de la salud, el coste se basará en el de un ingeniero biomédico ya que es la profesión que más se asemeja a este grado. El sueldo mensual aproximado de un ingeniero biomédico recién graduado ronda alrededor de los 1500 a los 1900 euros brutos \cite{CualEsSalario}.
Dado que se ha trabajado en este proyecto 6 meses a media jornada, el coste del empleado se reduce a los 850 euros, dejando su sueldo semestral a 5.100 euros. A esta cantidad hay que aplicarle las deducciones de la IRPF y cotizaciones a la Seguridad Social \cite{sanchez-toledoCualEsDiferencia2024}.

Las retenciones del IRPF no son fijas y dependen no solo de tu salario, sino que también de tus circunstancias personales y familiares. En este caso lo hemos establecido a 10\% ya que el empleado no tiene ni hijos a cargo ni discapacidades. Se puede observar en la tabla \ref{tab:salario_neto} cómo se queda el sueldo tras las deducciones correspondientes.

\begin{table}[h]
\centering
\begin{tabular}{lrr}
\hline
\textbf{Concepto}                  & \textbf{Porcentaje} & \textbf{Cuota (€)} \\ 
\hline
Contingencias comunes              & 4,70 \%             & 39,95              \\
Desempleo                          & 1,55 \%             & 13,18              \\
Formación Profesional              & 0,10 \%             & 0,85               \\
MEI                                & 0,12 \%             & 1,02               \\
\hline
\textbf{Total Seguridad Social}    &                     & 55,00              \\
\hline
Retención IRPF                     & 10,00 \%            & 85,00              \\
\hline
\textbf{Salario bruto}             &                     & 850,00             \\
\textbf{Salario neto}              &                     & 710,00             \\
\hline
\end{tabular}
\caption{Cálculo del sueldo neto del personal}
\label{tab:salario_neto}
\end{table}


\textbf{
\subsection{{Coste técnico}}}
\subsubsection{Costes de Software}

Todo el software empleado es de acceso abierto, por lo que no supuso ningún costo.

\subsubsection{Costes de Hardware}
En el caso de equipos externos, no se ha necesitado más que un ordenador portátil de la marca Lenovo, modelo exacto Lenovo Ideapad Slim 3 con un precio estimado de 624,49 euros, al que se le estima hasta 7  años de vida con buen mantenimiento y cuidado.
La amortización \footnote{Es la pérdida de valor de un bien o propiedad a lo largo del tiempo} del portátil durante el trabajo se ha calculado empleando la fórmula de \ref{fig:amortizacion_anual}.
\begin{figure}[h]
  \centering
  \[
    \text{Amortización anual}
    = \frac{\text{Valor de compra}}{\text{Vida útil estimada}}
  \]
  \caption{Fórmula de cálculo de la amortización anual}
  \label{fig:amortizacion_anual}
\end{figure}

Siendo la vida útil 7 años y el valor de 624,49 euros, la amortización anual se queda aproximadamente en 89,21 euros/año, en seis meses se queda la amortización de 45 euros aproximadamente.


\section{Viabilidad legal}
Al llevar a cabo este proyecto tecnológico, resulta imprescindible evaluar su viabilidad legal, no solo para proteger tanto el proceso como los derechos de propiedad intelectual, sino también para asegurar el cumplimiento de la normativa vigente actual.


\subsection{Datos y Consentimiento}

Esta aplicación no trata datos personales en ningún momento, por lo que su viabilidad legal está garantizada, ya que no se incumple ninguna normativa en materia de protección de datos. La plataforma ha sido desarrollada como una herramienta interactiva con fines exclusivamente educativos, divulgativos y preventivos en el ámbito de la salud pública, sin pretender en ningún caso sustituir la asistencia médica.
En primer lugar, la aplicación no recopila ni almacena datos personales de los usuarios. Las únicas interacciones previstas se limitan a la selección de parámetros no identificativos, como la zona geográfica o el fototipo de piel del usuario, lo cual la excluye del ámbito de aplicación del Reglamento General de Protección de Datos (RGPD).
En segundo lugar, todos los datos utilizados provienen de fuentes oficiales y de acceso público, concretamente de la Agencia Estatal de Meteorología (AEMET) y del Instituto Nacional de Estadística (INE). La legislación española permite el uso de estos datos con fines científicos, educativos o divulgativos, siempre que se respete la atribución de la fuente.
Finalmente, al tratarse de un Trabajo de Fin de Grado, queda justificado su carácter académico. Las herramientas implementadas en el resultado final consisten en un recomendador cuyas sugerencias se basan en condiciones ambientales, así como una estimación del riesgo acumulado, con el objetivo de alertar y fomentar el cuidado de la piel.

\subsection{Propiedad intelectual}
Respecto a la propiedad intelectual, el desarrollo de esta aplicación conlleva a la generación de una obra original que queda protegida por el Real Decreto legislativo 1/1996 de 12 de abril \cite{BOEA19968930RealDecreto}.Esta normativa establece los derechos que corresponden al autor sobre sus creaciones originales incluyendo programas informáticos.
Cabe destacar que al tratarse de un Trabajo de Fin de Grado no tendrá fines comerciales y por ello este proyecto se encuentra en un repositorio público.

\subsubsection{Logo de la aplicacion}
El logo de la aplicación cumple las leyes de propiedad intelectual debido a que es de elaboración propia con la ayuda de la inteligencia artificial. 